\section*{Chapitre \ref{ch:g5:irreversible-changes} --- Les transformations sans retour}

\begin{solution}{exo:g5:irreversible-changes:1}
Faire fondre du chocolat~: peut être défait (il durcit de nouveau en
refroidissant). \omterm{def:g4:what-a-fire-needs:burning}{Brûler} une allumette~: non. Dissoudre du sel~: peut être
défait (on fait évaporer l'eau). Griller du pain~: non.
\end{solution}

\begin{solution}{exo:g5:irreversible-changes:2}
Une transformation au cours de laquelle apparaissent de nouvelles
substances. Dans la cuisine~: faire cuire un œuf, faire cuire un gâteau
(ou encore~: griller du pain, une pomme coupée qui brunit).
\end{solution}

\begin{solution}{exo:g5:irreversible-changes:3}
En faisant évaporer l'eau~: le sel reste dans la coupelle. La
dissolution est une \omterm{def:g5:irreversible-changes:reversible}{transformation réversible}.
\end{solution}

\begin{solution}{exo:g5:irreversible-changes:4}
Une nouvelle odeur, et des grumeaux (un solide) qui apparaissent dans
le lait.
\end{solution}

\begin{solution}{exo:g5:irreversible-changes:5}
Les transformations de gauche~: le sucre qui \omterm{def:g2:mixing-and-dissolving:dissolve}{se dissout} dans l'eau, la
glace qui fond, le sable et le gravier mélangés. La double flèche
signifie que la transformation peut aller dans les deux sens~: on peut
la défaire.
\end{solution}

\begin{solution}{exo:g5:irreversible-changes:6}
L'huile protège le fer de la chaîne de l'\omterm{def:g4:air-a-mixture-of-gases:air}{air} et de l'eau~: elle ralentit
la rouille.
\end{solution}

\begin{solution}{exo:g5:irreversible-changes:7}
Des bulles de gaz qui apparaissent dans un liquide sans qu'on le
chauffe.
\end{solution}

\begin{solution}{exo:g5:irreversible-changes:8}
Utiles~: cuire les aliments, faire du pain, \omterm{def:g4:what-a-fire-needs:burning}{brûler} du gaz pour se
chauffer. Nuisibles~: la rouille, les aliments qui s'abîment, le bois
qui pourrit.
\end{solution}

\begin{solution}{exo:g5:irreversible-changes:9}
Le froid ralentit les \omterm{def:g5:irreversible-changes:chemical-change}{transformations chimiques} qui font tourner le
lait~: il se conserve plus longtemps.
\end{solution}

\begin{solution}{exo:g5:irreversible-changes:10}
Non. Les bulles sont de la vapeur d'eau, qui redevient de l'eau en
refroidissant~: rien de nouveau n'est fabriqué, et la transformation
peut être défaite. Des bulles seules ne prouvent pas qu'une nouvelle
substance est apparue.
\end{solution}

\begin{solution}{exo:g5:irreversible-changes:11}
Une bougie qui \omterm{def:g4:what-a-fire-needs:burning}{brûle} dégage de la chaleur et de la lumière, et fabrique
de nouvelles substances~: de la vapeur d'eau (un verre froid se couvre
de buée au-dessus d'elle) et du dioxyde de carbone. La cire qui \omterm{def:g4:what-a-fire-needs:burning}{brûle} ne
revient pas. La cire fondue, elle, durcit de nouveau en refroidissant.
\end{solution}
